\section{Threats to Validity}
\label{sec:threats}

We follow the classification of Wohlin et al.~\cite{wohlin2024experimentation} and list only threats specific to this study~\cite{verdecchia2023threats}.

\subsection{Construct validity}
\textbf{Executed is not able to detect.} RQ2 counts whether a benchmark executes a changed function~\cite{sjoberg2022construct}. Execution is necessary for detection, not sufficient: a benchmark may run the function on a trivial path. RQ2 is therefore an upper bound on what the suite can detect, and RQ3 measures the power side separately. The bias of RQ2 is towards \emph{fewer} undetectable edits than there are.
\textbf{What counts as an optimization edit.} The edit set depends on labels. All labels come from LLM agents: two blind raters and an adjudicator, all of the same model family; no human labelled the rows. Their agreement ($\kappa=\LabelKappa$) therefore overstates independent agreement, because raters of one model share its errors. One adjudicated label (a geopandas commit) cited details that the labelling sheet's truncated diff did not show; checked against the full commit diff, most cited details are present and one is not, and the label stands. The unreached count changes by one edit between the agent labels and an earlier reading of the commit subjects alone, and the interval is reported for both; neither set is human-validated.
\textbf{Synthetic slowdowns.} Calling a function body twice is a controlled probe of detectability, not a model of real regressions; it usually produces a large slowdown in the slowed function, so detection rates of moderate real regressions are likely lower than ours.
\textbf{Profiles at one state, one call.} Reach was profiled with the $S_0$ suite at $S_0$ for all edits, one call per case, not at each edit's parent. A lookup by qualified name found \ConstructUnreached{} edits unreached by construction (Section~\ref{sec:rq2}), but a benchmark added upstream between $S_0$ and an edit would be missed, which biases the unreached count upward. A re-profile reproduced the reach sets it covered, but profiles are samples of reach (Section~\ref{sec:lessons}).
\textbf{Time shares are approximate.} The share of a case's time spent in a slowed function comes from cProfile, whose overhead inflates call-heavy code; we use it only to classify misses, where the shares are far from the decision threshold.

\subsection{Internal validity}
\textbf{Machine noise and drift.} All timings come from one machine under WSL2, where the governor cannot be pinned; per-process setup cost drifted during the day for unknown reasons. Verdicts compare interleaved samples, which shields ratios from slow drift but not from changes in variance; cost is reported as ranges over passes. Some runs exceeded the load gate and are labelled noisy; none was silently dropped. \TODO{Drift passes with telemetry (E6) and, if granted, a pinned native-Linux machine for E2, E3 and E5.}
\textbf{Post-hoc analyses and deviations.} Several analyses were added after seeing data, and several protocol steps were not run as pre-registered; they are listed in Section~\ref{sec:deviations} and labelled where used. The primary pre-registered results (regression safety not measurable; the profile-completeness flag for \#9; the accumulation null) are reported unchanged.
\textbf{Concurrent work.} Other jobs on the same machine overlapped some runs; overlaps are logged per run, and those runs carry a noisy or concurrent label.
\textbf{Clone directory confounded with the edit.} Every A/B comparison ran the unpatched side in one clone directory and the patched side in another, so a verdict cannot separate the patch from differences between the clones; a repeat found off-target flips of both signs on about \TwoFlipPct\pct{} of cases. This biases off-target verdicts towards false flags; the on-target detection of \#2, which repeats at a ratio far above its noise floor, is unlikely to be explained by it. Swapping or rotating the directories between sides would remove the confound.

\subsection{External validity}
\textbf{One project for the main results.} RQ2 and RQ3 rest on networkx, with \NxAgentOpt{} optimization edits; the coverage interval is wide (\ci{\BOneAgentLo}{\BOneAgentHi}). Our static pilot on other projects is not evidence (Appendix~\ref{app:q3}). Every result is stated for networkx until the replication lands. \TODO{E1 replication on 3--4 projects; per-project and pooled results.}
\textbf{Replayed commits stand in for a session.} Historical commits are larger and further apart than an optimizer's candidates. Per-edit results (coverage, power) do not depend on the session model; cost totals and the accumulation policy do.
\textbf{Python, asv, small builds.} The build results hold for pure-Python and small C-extension builds; we did not study a project with a heavy build.

\subsection{Conclusion validity}
\textbf{Small samples.} Proportions over \NxAgentOpt{} edits, \NumSlowdowns{} slowdowns and \MutantsKilled{} mutants have wide intervals, which we report everywhere; ``no miss'' means a miss rate below the interval's upper bound, not safety.
\textbf{Multiplicity.} The oracle made \OffTargetVerdictable{} off-target comparisons that can yield a verdict at $\alpha=\AsvAlpha$; about \OffTargetExpected{} false verdicts would be expected at the nominal level and \OffTargetObserved{} were observed, but the null rate of the joint rule is not calibrated here, and samples from one process are not independent, which makes Mann-Whitney $p$ optimistic. A repeat of \TwoTotal{} off-target comparisons produced \TwoFlips{} new non-``same'' verdicts (about \TwoFlipPct\pct{} in this fixed-directory design, source unresolved), so off-target and small on-target verdicts such as \#8 cannot be read individually. Both original off-target flags were first judged noise by a reach argument and did not repeat.
\textbf{Power of the A/A control.} \AAComparisons{} comparisons of \AACases{} cases cannot separate a small false-verdict rate from zero, and the heavy-module estimate rests on \AAAlgoCases{} cases. \TODO{E2 sigma.}
\textbf{Exact test.} Verdicts were first computed with a normal approximation; re-testing all of them exactly changed none.

\subsection{Reliability and research process}
\label{sec:deviations}
\textbf{Artifacts.} Frozen inputs carry hashes; raw files record the state and harness version; the paper's numbers and figures are generated from the raw files by a script that fails if they disagree. One raw file (the A/A $\sigma$ estimates) lacks its metadata block. \TODO{Artifact package and DOI.}
\textbf{Deviations from the pre-registration.} (1) The noise floor $1+3\sigma$ was applied to every flagged verdict, with $\sigma$ from five A/A repeats at each verdict's own state rather than at $S_0$ as planned; the other benchmark cases still have no $\sigma$. (2) The first out-of-reach spot check had no power and was replaced by a powered profile-completeness check. (3) The labelling rule for noisy test runs was written after the first look at two slow states. (4) Edit labels were added after mining, and the mining rule was not changed for this study. (5) A per-arm regression simulation was not run. \TODO{Full deviation table in the artifact; update after E2.}
\textbf{AI assistance.} The selection and validation rules contain no learned or LLM component; the commit labels that define the RQ2 edit set were produced by LLM agents (construct validity, above). Large parts of the harness, the analysis scripts and the text were also produced with AI assistance. \TODO{Disclose per the chosen venue's policy; the authors re-derive every headline number from the raw files before submission.}
